% Options for packages loaded elsewhere
\PassOptionsToPackage{unicode}{hyperref}
\PassOptionsToPackage{hyphens}{url}
\PassOptionsToPackage{dvipsnames,svgnames,x11names}{xcolor}
\documentclass[
  11pt,
  a4paper,
]{article}
\usepackage{xcolor}
\usepackage[margin=24mm]{geometry}
\usepackage{amsmath,amssymb}
\setcounter{secnumdepth}{5}
\usepackage{iftex}
\ifPDFTeX
  \usepackage[T1]{fontenc}
  \usepackage[utf8]{inputenc}
  \usepackage{textcomp} % provide euro and other symbols
\else % if luatex or xetex
  \usepackage{unicode-math} % this also loads fontspec
  \defaultfontfeatures{Scale=MatchLowercase}
  \defaultfontfeatures[\rmfamily]{Ligatures=TeX,Scale=1}
\fi
\usepackage{lmodern}
\ifPDFTeX\else
  % xetex/luatex font selection
\fi
% Use upquote if available, for straight quotes in verbatim environments
\IfFileExists{upquote.sty}{\usepackage{upquote}}{}
\IfFileExists{microtype.sty}{% use microtype if available
  \usepackage[]{microtype}
  \UseMicrotypeSet[protrusion]{basicmath} % disable protrusion for tt fonts
}{}
\makeatletter
\@ifundefined{KOMAClassName}{% if non-KOMA class
  \IfFileExists{parskip.sty}{%
    \usepackage{parskip}
  }{% else
    \setlength{\parindent}{0pt}
    \setlength{\parskip}{6pt plus 2pt minus 1pt}}
}{% if KOMA class
  \KOMAoptions{parskip=half}}
\makeatother
\usepackage{longtable,booktabs,array}
\newcounter{none} % for unnumbered tables
\usepackage{calc} % for calculating minipage widths
% Correct order of tables after \paragraph or \subparagraph
\usepackage{etoolbox}
\makeatletter
\patchcmd\longtable{\par}{\if@noskipsec\mbox{}\fi\par}{}{}
\makeatother
% Allow footnotes in longtable head/foot
\IfFileExists{footnotehyper.sty}{\usepackage{footnotehyper}}{\usepackage{footnote}}
\makesavenoteenv{longtable}
\setlength{\emergencystretch}{3em} % prevent overfull lines
\providecommand{\tightlist}{%
  \setlength{\itemsep}{0pt}\setlength{\parskip}{0pt}}
% Shared style; used by generated, standalone LaTeX sources.
\usepackage{xcolor}
\usepackage{xurl}
\usepackage{fvextra}
\usepackage{etoolbox}
\usepackage{fancyhdr}
\usepackage{titlesec}
\definecolor{researchblue}{HTML}{173B58}
\AtBeginDocument{\hypersetup{linkcolor=researchblue,urlcolor=researchblue}}
\pagestyle{fancy}
\fancyhf{}
\fancyhead[L]{\small\sffamily SVG PATCH LAB}
\fancyhead[R]{\small\sffamily CVPR 2027 research package}
\fancyfoot[L]{\small Working documents; results and proposals distinguished}
\fancyfoot[R]{\thepage}
\setlength{\headheight}{15pt}
\setlength{\emergencystretch}{3em}
\titleformat{\section}{\Large\bfseries\color{researchblue}}{\thesection}{0.7em}{}
\titleformat{\subsection}{\large\bfseries\color{researchblue}}{\thesubsection}{0.7em}{}
\DefineVerbatimEnvironment{verbatim}{Verbatim}{breaklines=true,breakanywhere=true,fontsize=\small}
\AtBeginEnvironment{longtable}{\small\setlength{\tabcolsep}{4pt}}
\let\originaltableofcontents\tableofcontents
\renewcommand{\tableofcontents}{\originaltableofcontents\clearpage}
\usepackage{bookmark}
\IfFileExists{xurl.sty}{\usepackage{xurl}}{} % add URL line breaks if available
\urlstyle{same}
\hypersetup{
  pdftitle={Experiment Protocol and Work Plan},
  pdfauthor={SVG Patch Lab - Research Working Package},
  colorlinks=true,
  linkcolor={Maroon},
  filecolor={Maroon},
  citecolor={Blue},
  urlcolor={blue},
  pdfcreator={LaTeX via pandoc}}

\title{Experiment Protocol and Work Plan}
\author{SVG Patch Lab - Research Working Package}
\date{14 September 2026 \textbar{} CVPR 2027 target}

\begin{document}
\maketitle

{
\setcounter{tocdepth}{2}
\tableofcontents
}
\section{Status and rules of
interpretation}\label{status-and-rules-of-interpretation}

This document separates completed diagnostics, runnable reproductions
and designed future experiments. A planned experiment has no result. A
working script is not evidence that a model was trained or that a
comparison was completed. The machine-readable registry is
\texttt{experiments/registry.json};
\texttt{scripts/experiments.py\ list} shows statuses and commands.

The main scientific question concerns numerical and visible semantic
correctness after engineering-drawing edits. The existing Qwen
compact-DOM pilot is a separate infrastructure check. Every main result
must identify the physical case, request, model/configuration, solver
reference, artifact and evaluator version.

\section{Completed experiments}\label{completed-experiments}

{\def\LTcaptype{none} % do not increment counter
\begin{longtable}[]{@{}
  >{\raggedright\arraybackslash}p{(\linewidth - 6\tabcolsep) * \real{0.2500}}
  >{\raggedright\arraybackslash}p{(\linewidth - 6\tabcolsep) * \real{0.2500}}
  >{\raggedright\arraybackslash}p{(\linewidth - 6\tabcolsep) * \real{0.2500}}
  >{\raggedright\arraybackslash}p{(\linewidth - 6\tabcolsep) * \real{0.2500}}@{}}
\toprule\noalign{}
\begin{minipage}[b]{\linewidth}\raggedright
ID
\end{minipage} & \begin{minipage}[b]{\linewidth}\raggedright
Scope
\end{minipage} & \begin{minipage}[b]{\linewidth}\raggedright
Evidence
\end{minipage} & \begin{minipage}[b]{\linewidth}\raggedright
Important limit
\end{minipage} \\
\midrule\noalign{}
\endhead
\bottomrule\noalign{}
\endlastfoot
C01 & Evaluator and edit-policy regressions & 30 passing tests & Tests
cover implemented scope, not all SVG semantics. \\
C02 & Manufactured contour corruptions & 160 controls; numerical audit
accepts 20/120 corruptions & All remaining false accepts are wrong
labels. \\
C03 & Historical v3 re-audit & Six attempts, three strict geometric
passes & Corner and mask effects prevent blanket physics-failure
claims. \\
C04 & Rule baseline & 60/60 test and 60/60 annulus & Shared templates
make this pilot trivial. \\
C05 & A100 LoRA training & Three seeds, three epochs, 90 steps each;
final adapters saved & Same perfect scores as the rule baseline. \\
C06 & Astra API recheck & Four completed calls with raw responses and
token usage & One sample per condition; not a failure-rate estimate. \\
C07 & Independent Robin FEM & Five mesh resolutions, analytic anchor,
heat balance & Mesh differences are estimates, not certified error
bounds. \\
C08 & Direct FEM contour scoring & Maxima 0.111 and 0.374 degrees
Celsius & Finite geometry sampling does not certify the entire
drawing. \\
C09 & Historical capacitor path audit & 12 open field paths; no alleged
closed loops & Source-specific geometry check, not electrostatic
validation. \\
C10 & Upstream FEM-Bench SVG bridge & Two upstream tests and 24 FEM
solves & A one-dimensional integration anchor only. \\
\end{longtable}
}

Rerunning CPU reproductions should write to a new directory. Preserve
original measured outputs and manifests. API or GPU repeats require an
explicit job command; the default runner never launches them
automatically.

\section{P01: numerical reference
suite}\label{p01-numerical-reference-suite}

\textbf{Purpose:} establish trustworthy references before model scoring.
Status: designed; Robin anchor implemented, other families incomplete.

Heat: steady two-dimensional conduction on rectangles and nonconvex
domains with Dirichlet, Neumann and Robin boundaries. Record
conductivity and length units. Validate Robin signs and corner treatment
using an analytic manufactured case and boundary-flux balance. Use at
least three mesh resolutions; choose the final resolution based on field
and relevant contour-point convergence.

Elasticity: finite plates with holes and rounded brackets, with explicit
geometry, Young's modulus, Poisson ratio, plane-stress/plane-strain
choice, support model and distributed traction. Check displacement,
force balance, free-surface traction and stress recovery. Treat
sharp-corner peak stress as unresolved unless the problem explicitly
asks for a singularity/disclosure response. Use the exact finite
geometry rather than substituting a remote infinite-domain formula.

Electrostatics: specify 2-D cross-section or 3-D conductor geometry,
thickness, potential values and the exterior boundary model. Perform
both mesh and exterior-domain refinement. Check conductor potential,
symmetry, charge/flux balance and orthogonality of field lines and
equipotentials. An image that ends its contours inside the canvas does
not define behavior at infinity.

\textbf{Acceptance:} reference metadata, physical-balance checks,
convergence tables and numerical-error estimates accompany every case.
Cases with unresolved reference uncertainty remain in a separate
diagnostic set and do not become hidden failures or exclusions.

\section{P02: claim schema and complete
verifier}\label{p02-claim-schema-and-complete-verifier}

\textbf{Purpose:} bind visible content to physical quantities. Status:
partial geometric implementation; semantic and occlusion coverage
planned.

Specify solution identifiers, quantity names, units, contour levels,
coordinate maps, label anchors and legend mappings. Define invalidation
for geometry, material, loads and boundary-condition changes. Add tests
for ancestor transforms, wrong scalar/tensor quantities, altered units,
stale solution records, missing contours, hidden layers, clipping and
label/legend mismatch.

Measure false acceptance on both malicious and accidental corruptions.
Include honest edits that should be permitted, such as changing an
incorrect label to the correct value. Do not rely on a string hash as
evidence of visible equivalence. Unsupported features receive an
explicit status rather than a silent pass.

\textbf{Acceptance:} the checker rejects its specified corruptions and
accepts matched honest counterfactuals; unsupported features and
remaining blind spots are documented. A full-drawing certificate must
not be reported until its visible-semantic scope is actually implemented
and validated.

\section{P03: main case and edit
collection}\label{p03-main-case-and-edit-collection}

\textbf{Purpose:} avoid a shared-template benchmark that a rule parser
solves. Status: designed.

Proposed initial collection: 240 physical cases across three families,
split before generating edit variants into 120 training, 48 validation
and 72 test cases. These are planning counts and may change after a
reference-only feasibility pilot; record any change before model
comparisons. Store separate geometry-family and instruction-template
holdouts rather than calling every new parameter value
out-of-distribution.

Each case should include presentation edits, label/legend edits, a
physical-input edit, an under-specified request and an honest/misleading
counterfactual pair where appropriate. Use independent human-written
language or externally sourced editing requests for part of evaluation.
Retain one instruction provenance record per request. Do not train on
the test references' prompts, outputs or inspected failure corrections.

\textbf{Acceptance:} no physical-case leakage; reproducible seeds;
reference eligibility fixed; editing instructions require context beyond
keyword routing; task counts and exclusions are recorded before
evaluation.

\section{P04: primary model and deterministic
comparisons}\label{p04-primary-model-and-deterministic-comparisons}

\textbf{Purpose:} test whether the proposed representation helps at
equal information. Status: designed.

{\def\LTcaptype{none} % do not increment counter
\begin{longtable}[]{@{}
  >{\raggedright\arraybackslash}p{(\linewidth - 4\tabcolsep) * \real{0.3333}}
  >{\raggedright\arraybackslash}p{(\linewidth - 4\tabcolsep) * \real{0.3333}}
  >{\raggedright\arraybackslash}p{(\linewidth - 4\tabcolsep) * \real{0.3333}}@{}}
\toprule\noalign{}
\begin{minipage}[b]{\linewidth}\raggedright
Baseline
\end{minipage} & \begin{minipage}[b]{\linewidth}\raggedright
Input and freedom
\end{minipage} & \begin{minipage}[b]{\linewidth}\raggedright
Main comparison
\end{minipage} \\
\midrule\noalign{}
\endhead
\bottomrule\noalign{}
\endlastfoot
Direct Astra & Physical task and drawing; generates/edits SVG & Unaided
final-artifact fidelity. \\
Solver-informed unrestricted editor & Same numerical solution supplied;
unrestricted SVG edits & Value of enforced protection beyond solver
access. \\
Published solver-agent workflow & Pinned runnable
MechAgents/MCP-SIM/OASiS adaptation & End-to-end engineering workflow
comparison. \\
Protected geometry only & Numerical path protection; weaker annotation
binding & Contribution of semantic checking. \\
Joint geometry and semantics & Proposed representation and final
verifier & Full proposed method. \\
Deterministic solver plot / template & Solver-owned geometry and labels
& Necessity and usefulness of learned drafting. \\
Rigid freeze baseline & Minimal permitted style changes & Safety versus
editing flexibility. \\
\end{longtable}
}

Use identical physical specifications, references and requests where
applicable. Label differences in solver access explicitly. Preserve all
attempted outputs, including invalid SVG and API failures. API failures
are operational outcomes, not evidence of PDE incapability. Record model
revision/alias, effort, token cap, endpoint, wall time and usage.

\textbf{Acceptance:} per-case results for every baseline and the
complete attempted denominator. No aggregate success claim from only
successfully rendered or numerically supported cases.

\section{P05: training study}\label{p05-training-study}

\textbf{Purpose:} test a learned improvement on a task that needs it.
Status: designed; the earlier template pilot is complete but does not
answer this question.

Start only after P03/P04 show a gap beyond the rule and deterministic
controls. Train the base model and final SFT policy with the same
executor and information. Record the source, filtering and
label-generation process. Use completion-only supervision for structured
actions when appropriate. If the main claim is multimodal, select a
model that genuinely sees the image or complete graphical
representation; do not relabel a compact DOM index as vision input.

Choose a small validation-driven hyperparameter search before a
three-seed final configuration. Record every tried configuration. Keep
the evaluation set fixed. Report strict output validity, instruction
success and physical correctness separately. Compare base,
prompting-only, SFT and deterministic policies. Reuse A100 for
compatible small-model runs; choose larger hardware only after measuring
memory and throughput needs.

\textbf{Acceptance:} a reproducible, useful held-out gain; or a
documented negative result. Perfect performance on shared templates is
not a trigger to scale training.

\section{P06: sequential physical and visual
edits}\label{p06-sequential-physical-and-visual-edits}

\textbf{Purpose:} expose stale fields and broken bindings across turns.
Status: designed.

Create prespecified sequences of five edits: presentation change,
physical parameter change, annotation rearrangement, geometry or
boundary change, and a request that could reuse an old claim
incorrectly. Include shorter two-edit anchors for debugging. Every
physical change must invoke a real solver and produce an updated
reference record; rerouting alone does not count.

Track per-step success, first failure, stale-solution acceptance and
final visible correctness. Also track cumulative latency and solver
calls. Reusing cached solutions is permitted only if the physical
specification is unchanged and the cache key identifies the complete
problem.

\textbf{Acceptance:} complete sequence traces, actual solver evidence,
and results against a deterministic replotting baseline.

\section{P07: ablations and
robustness}\label{p07-ablations-and-robustness}

\textbf{Purpose:} identify which components explain any improvement.
Status: designed.

Remove, one at a time, the solver record, protected geometry, unit
binding, legend binding, visibility checks and post-edit verification.
Compare path-string freezing to transformed-geometry checking. Evaluate
instruction paraphrases, changed element IDs, nested transforms,
alternate layout templates and previously unseen geometry families.
Include beneficial semantic corrections as well as deceptive edits.

Do not retrofit corruptions solely around known strengths of the
proposed verifier. Freeze a public development set and a separately
constructed evaluation set. Explain each ablation's changed capabilities
so that losing useful edit freedom cannot masquerade as a better safety
method.

\section{P08: drafting and visible correctness
assessment}\label{p08-drafting-and-visible-correctness-assessment}

\textbf{Purpose:} establish visual usefulness for a CVPR submission.
Status: designed.

Prepare anonymized, randomized outputs from matched cases. Ask raters to
judge instruction fulfillment, readability, numerical-label consistency
and suitability for the specified drafting task. Use three independent
raters where feasible, report agreement, and include the exact rubric.
Choose the assessed sample before viewing preferred-model outcomes.
Report numerical audits separately from visual ratings; an attractive
drawing can still be physically wrong.

Use a pilot to estimate time and reliability before fixing the final
sample size. The current package contains no completed human study or
fabricated ratings.

\section{Metrics and uncertainty}\label{metrics-and-uncertainty}

Report curve-field mean and maximum error normalized by an explicitly
specified field range, required-level coverage, unsupported samples,
quantity/unit/legend errors, edit fulfillment, render validity,
rejection/disclosure correctness, solver calls, latency and tokens.
State whether errors come from direct FEM evaluation or regular-grid
interpolation.

Use physical cases as the resampling unit. Related edits from one case
and repeated seeds on the same test set are not independent cases.
Report paired case-level differences and confidence intervals when
sample size permits. For sequences, resample complete sequences or their
source cases. Publish raw values alongside uncertainty summaries.

Separate reference error, output error and operational failure. Report
both full-domain metrics and any prespecified singular-region analysis.
A posteriori exclusions are diagnostics and must not silently replace
the main score.

\section{Compute and stopping rules}\label{compute-and-stopping-rules}

Begin with a small end-to-end feasibility batch. Estimate tokens per
case and wall time from that batch before scheduling a large API study.
Record a maximum job count and budget in the run configuration. Training
jobs should checkpoint and export automatically; the recovered A100 runs
demonstrate why local artifact retention is necessary.

Stop or redesign if references are unresolved, a task is
under-specified, deterministic controls solve the main task, or the
proposed verifier achieves safety only by refusing useful edits. Use
validation for model selection. Do not keep sampling until a preferred
outcome appears.

\section{Submission deliverables and
dependencies}\label{submission-deliverables-and-dependencies}

P01/P02 precede main scoring. P03 precedes P04/P05. P06/P07/P08 support
the final contribution claim. A final package needs complete results,
limitations, verified bibliography, figure sources, anonymous
paper/supplement, licenses and author review. Official registration,
paper and supplement dates should be rechecked before submission.

Current status: the documents and experimental infrastructure are
packaged; the main benchmark, stronger learned study, full semantic
verifier and human assessment remain to be executed. The working paper
is explicitly unfinished.

\end{document}
